% =============================================================================
%  Reproducibility Report – Coral Reef Terrain Analysis (USVI)
%  Compile with: pdflatex REPRODUCIBILITY_REPORT.tex  (twice for TOC/refs)
%  Required packages: geometry, hyperref, booktabs, longtable, xcolor,
%                     listings, fancyhdr, titlesec, parskip, microtype
% =============================================================================

\documentclass[11pt, a4paper]{article}

% ── Encoding & fonts ──────────────────────────────────────────────────────────
\usepackage[T1]{fontenc}
\usepackage[utf8]{inputenc}
\usepackage{lmodern}
\usepackage{microtype}

% ── Page geometry ─────────────────────────────────────────────────────────────
\usepackage[top=2.5cm, bottom=2.5cm, left=2.8cm, right=2.8cm]{geometry}

% ── Colour palette ────────────────────────────────────────────────────────────
\usepackage[dvipsnames, table]{xcolor}
\definecolor{coralblue}{RGB}{0,84,147}      % main heading colour
\definecolor{coralred}{RGB}{200,60,40}      % accent / warnings
\definecolor{lightgray}{RGB}{245,245,245}   % code background
\definecolor{midgray}{RGB}{200,200,200}     % table rules
\definecolor{darkgray}{RGB}{80,80,80}       % secondary text

% ── Hyperlinks ────────────────────────────────────────────────────────────────
\usepackage[colorlinks=true,
            linkcolor=coralblue,
            urlcolor=coralblue,
            citecolor=coralblue,
            pdftitle={Reproducibility Report – Coral Terrain Analysis USVI},
            pdfauthor={J. Olaya, GitHub Copilot Coding Agent},
            pdfsubject={Spatial ecology – benthic terrain predictors},
            pdfkeywords={coral reef, USVI, terrain analysis, ArcGIS, reproducibility}
           ]{hyperref}

% ── Section styling ───────────────────────────────────────────────────────────
\usepackage{titlesec}
\titleformat{\section}{\large\bfseries\color{coralblue}}{}{0em}{}[\color{midgray}\titlerule]
\titleformat{\subsection}{\normalsize\bfseries\color{coralblue}}{}{0em}{}
\titleformat{\subsubsection}{\normalsize\itshape\color{darkgray}}{}{0em}{}

% ── Header / footer ───────────────────────────────────────────────────────────
\usepackage{fancyhdr}
\pagestyle{fancy}
\fancyhf{}
\fancyhead[L]{\small\color{darkgray}Reproducibility Report — Coral Terrain Analysis, USVI}
\fancyhead[R]{\small\color{darkgray}April 2026}
\fancyfoot[C]{\small\thepage}
\renewcommand{\headrulewidth}{0.4pt}

% ── Paragraph spacing ─────────────────────────────────────────────────────────
\usepackage{parskip}
\setlength{\parindent}{0pt}
\setlength{\parskip}{6pt}

% ── Tables ────────────────────────────────────────────────────────────────────
\usepackage{booktabs}
\usepackage{longtable}
\usepackage{array}
\usepackage{tabularx}
\newcolumntype{L}[1]{>{\raggedright\arraybackslash}p{#1}}
\newcolumntype{C}[1]{>{\centering\arraybackslash}p{#1}}

% ── Code listings ─────────────────────────────────────────────────────────────
\usepackage{listings}
\lstset{
  backgroundcolor=\color{lightgray},
  basicstyle=\ttfamily\small,
  breaklines=true,
  frame=single,
  framerule=0.4pt,
  rulecolor=\color{midgray},
  showstringspaces=false,
  columns=flexible,
}
\lstdefinestyle{Rcode}{language=R, keywordstyle=\color{coralblue}\bfseries,
                        stringstyle=\color{coralred}}
\lstdefinestyle{pycode}{language=Python, keywordstyle=\color{coralblue}\bfseries,
                         stringstyle=\color{coralred}}
\lstdefinestyle{shell}{basicstyle=\ttfamily\footnotesize\color{darkgray},
                        backgroundcolor=\color{lightgray}}

% ── Miscellaneous ─────────────────────────────────────────────────────────────
\usepackage{graphicx}
\usepackage{enumitem}
\usepackage{amsmath}
\usepackage{amssymb}
\usepackage{pifont}   % \ding{51} = ✓
\newcommand{\cmark}{\textcolor{OliveGreen}{\ding{51}}}
\newcommand{\xmark}{\textcolor{coralred}{\ding{55}}}

% ── Date ──────────────────────────────────────────────────────────────────────
\usepackage{datetime2}

% =============================================================================
%  TITLE PAGE
% =============================================================================

\begin{document}

\begin{titlepage}
  \centering
  \vspace*{2cm}

  {\color{coralblue}\rule{\linewidth}{1.5pt}}\\[0.4cm]

  {\LARGE\bfseries\color{coralblue}
    Reproducibility Report\\[0.3em]
    Coral Reef Terrain Analysis\\[0.3em]
    St.\ Thomas / St.\ John, US Virgin Islands (USVI)
  }\\[0.4cm]

  {\large\color{darkgray}
    Spatial Predictor Stack Construction at 50\,m / EPSG:32620
  }\\[0.3cm]

  {\color{coralblue}\rule{\linewidth}{1.5pt}}\\[1.5cm]

  \begin{tabular}{rl}
    \textbf{Domain scientist:}  & J.\ Olaya \\[4pt]
    \textbf{AI assistant:}      & GitHub Copilot Coding Agent (Claude Sonnet) \\[4pt]
    \textbf{Date:}              & April 2026 \\[4pt]
    \textbf{Repository:}        & \href{https://github.com/jolaya80/jolaya80}%
                                       {\texttt{github.com/jolaya80/jolaya80}} \\[4pt]
    \textbf{Main script:}       & \texttt{Coral\_Modeling\_Terrain\_Analysis\_USVI\_202604.py} \\[4pt]
    \textbf{Supporting scripts:}& \texttt{01\_exploratory\_spatial\_analysis.R} \\
                                & \texttt{02\_reproject\_resample.R} \\
  \end{tabular}

  \vfill

  {\small\color{darkgray}
    \textit{This document is intended to enable full reproducibility of the
    terrain predictor pipeline by any researcher with access to ArcGIS Pro
    (Spatial Analyst Extension) and/or R\,($\geq$\,4.1).}
  }
\end{titlepage}

% =============================================================================
%  TABLE OF CONTENTS
% =============================================================================

\tableofcontents
\newpage

% =============================================================================
\section{Purpose of This Document}
% =============================================================================

This report documents the complete workflow used to derive bathymetric terrain
predictors for coral-cover modelling in St.\ Thomas and St.\ John (USVI).  Its
primary objective is \textbf{reproducibility}: any researcher with the required
software and data should be able to re-run every step and obtain identical
outputs.

The document records:

\begin{itemize}[noitemsep]
  \item the original research question and spatial specifications provided by
        the domain scientist;
  \item the technical contributions made by the AI coding agent;
  \item a chronological account of the human–AI interaction;
  \item all bugs found and corrected during development;
  \item the final verified outputs; and
  \item step-by-step instructions for non-expert users.
\end{itemize}

% =============================================================================
\section{Project Background and Initial Request}
% =============================================================================

\subsection{Research Context}

The work belongs to a coral-cover modelling project for St.\ Thomas and
St.\ John (USVI), funded under the NSF CoPE programme
(\textit{Coral Reefs in a Rapidly Changing Environment}).
The goal is to predict benthic hard-coral cover percentages from environmental
predictors using a statistical model trained on field-survey observations.

Terrain descriptors derived from high-resolution (2\,m) bathymetry serve as
the spatial predictors.  Table~\ref{tab:predictors} lists those chosen and
their ecological rationale.

\begin{table}[h]
\centering
\caption{Terrain predictors and ecological justification.}
\label{tab:predictors}
\rowcolors{2}{lightgray}{white}
\begin{tabularx}{\linewidth}{L{5cm}X}
\toprule
\textbf{Predictor} & \textbf{Ecological rationale} \\
\midrule
Slope (local, 50\,m)        & Substrate angle; affects sedimentation and larval settlement \\
Slope (broad-scale, 240\,m) & Regional exposure to currents and wave energy \\
Slope-of-slope              & Terrain ruggedness; structural complexity proxy \\
Aspect                      & Directional orientation; related to light and current exposure \\
Curvature                   & Concavity / convexity; controls sediment convergence \\
\bottomrule
\end{tabularx}
\end{table}

\subsection{The Initial Request}

The user's original task (paraphrased) was:

\begin{quote}
\textit{``I need to build spatial terrain predictors for St.\ Thomas USVI from
a 2\,m bathymetry raster.  The target coordinate reference system is WGS\,84 /
UTM Zone\,20N (EPSG:32620) and the final output resolution must be
50\,m $\times$ 50\,m.  I need slope at multiple scales, slope-of-slope,
aspect, and curvature.  Can you write a Python script using ArcGIS\,Pro /
\texttt{arcpy} to do this?''}
\end{quote}

The user shared a preliminary Python script and asked the AI to review, correct,
and extend it.  Subsequent messages involved iterative debugging as issues were
discovered during execution.

% =============================================================================
\section{User Contributions}
% =============================================================================

The domain scientist provided all of the context in Table~\ref{tab:user}.
Without this information the AI could not have constructed the correct analysis.

\begin{table}[h]
\centering
\caption{Knowledge and criteria provided by the domain scientist.}
\label{tab:user}
\rowcolors{2}{lightgray}{white}
\begin{tabularx}{\linewidth}{L{4.5cm}X}
\toprule
\textbf{Category} & \textbf{Detail} \\
\midrule
Study area            & St.\ Thomas + St.\ John (USVI) \\
Research programme    & NSF CoPE \\
Target predictors     & Slope (50\,m \& 240\,m), slope-of-slope, aspect, curvature \\
Source data           & \texttt{STTSTJ\_2m.tif} (2\,m bathymetry) on shared Google Drive \\
Target CRS            & EPSG:32620 — WGS\,84 / UTM Zone\,20N \\
Target resolution     & 50\,m $\times$ 50\,m \\
Focal neighbourhood   & 240\,m radius for broad-scale slope \\
Folder structure      & NSF CoPE internal Google Drive path \\
Tooling               & ArcGIS\,Pro + \texttt{arcpy} (Python); R for QA \\
QA criteria           & All outputs exactly 50.0\,m / EPSG:32620; keep 2\,m archives \\
Execution feedback    & Error messages and unexpected file sizes from each run \\
\bottomrule
\end{tabularx}
\end{table}

% =============================================================================
\section{AI Agent Contributions}
% =============================================================================

\subsection{Code Architecture}

\begin{itemize}[noitemsep]
  \item Block-by-block Jupyter Notebook structure (BLOCK\,0–8) allowing each
        step to be run and verified independently.
  \item \texttt{WorkflowConfig} class centralising all parameters (paths, CRS,
        resolutions, aggregation factors, focal radii).
  \item \texttt{Logger} utility for timestamped console output.
  \item \texttt{\_save\_raster()} helper working around silent save failures on
        Google Drive for large rasters.
\end{itemize}

\subsection{Technical Diagnosis and Bug Fixes}

Seven critical bugs were identified through analysis of the user's reported
errors and the initial code (see Section~\ref{sec:fixes} for full details).

\subsection{R Scripts}

\begin{itemize}[noitemsep]
  \item \texttt{01\_exploratory\_spatial\_analysis.R} — loads all raster and
        vector inputs, clips the WCMC coral-reef polygon mask to the study area,
        reports CRS/resolution, and counts survey points inside the coral mask.
  \item \texttt{02\_reproject\_resample.R} — reprojects and resamples all
        terrain rasters from EPSG:26920 to EPSG:32620 at 50\,m using R's
        \texttt{terra} package, including circular-mean aspect aggregation.
        (No longer required after the Python fixes; retained for reference.)
\end{itemize}

\subsection{Documentation}

All code is extensively commented, explaining \textit{what} each step does and
\textit{why} — ecological rationale, statistical correctness, and the risks of
alternative approaches.  Each BLOCK begins with a docstring listing the fix
applied, expected outputs, and estimated processing time.

% =============================================================================
\section{Conversation Flow and Collaborative Process}
% =============================================================================

\subsection{How the Interface Works}

The work was done using the \textbf{GitHub Copilot Coding Agent} (Task Agent),
which operates directly on a GitHub repository:

\begin{enumerate}[noitemsep]
  \item The user submits a task description in natural language via the
        GitHub Copilot interface.
  \item The agent clones the repository into a sandboxed environment, reads all
        existing files, and plans the work.
  \item The agent writes, edits, and commits code changes directly to a feature
        branch.
  \item The user reviews all changes via Pull Requests on GitHub.
  \item The agent reports progress at each meaningful step.
  \item Conversations are iterative: the user can submit follow-up tasks, error
        messages, or new requirements, and the agent continues on the same
        branch.
\end{enumerate}

\subsection{Chronological Summary}

Table~\ref{tab:conversation} summarises each stage of the dialogue.

\begin{longtable}{L{3cm}L{6cm}L{5.5cm}}
\caption{Summary of human–AI interaction.}\label{tab:conversation}\\
\toprule
\textbf{Stage} & \textbf{User action} & \textbf{AI action} \\
\midrule
\endfirsthead
\multicolumn{3}{c}{\small\textit{(continued from previous page)}}\\
\toprule
\textbf{Stage} & \textbf{User action} & \textbf{AI action} \\
\midrule
\endhead
\midrule
\multicolumn{3}{r}{\small\textit{(continued on next page)}}\\
\endfoot
\bottomrule
\endlastfoot
Initial request
  & Described research goal, data, and required outputs
  & Reviewed code; identified 7 bugs; proposed corrected architecture \\[4pt]
First code version
  & Ran BLOCK\,0 and BLOCK\,1; confirmed workspace set up
  & Committed Blocks 0–2 with fixes 1 \& 2 \\[4pt]
Slope computation
  & Ran BLOCK\,3; reported unexpected file sizes
  & Diagnosed wrong \texttt{AGG\_FACTOR}; fixed single-step aggregation \\[4pt]
240\,m slope
  & Ran BLOCK\,4; processing too fast ($<$10\,s, expected 5–8\,min)
  & Diagnosed wrong focal radius (5 cells instead of 120) \\[4pt]
Aspect \& curvature
  & Ran Blocks 6–7; asked about circular mean
  & Explained 0°/360° wrap-around; implemented sin/cos decomposition;
    added 50\,m curvature aggregation in Python \\[4pt]
Final verification
  & Ran BLOCK\,8; all five outputs \cmark{} 50\,m / EPSG:32620
  & Confirmed workflow complete; noted R script no longer needed \\[4pt]
Documentation
  & Requested reproducibility report
  & Wrote Markdown and \LaTeX{} versions \\
\end{longtable}

% =============================================================================
\section{Technical Fixes Applied}\label{sec:fixes}
% =============================================================================

Table~\ref{tab:fixes} documents all seven corrections.  Each fix is labelled in
the script header and in the relevant BLOCK docstring.

\begin{longtable}{C{0.6cm}L{2cm}L{5.2cm}L{5.2cm}}
\caption{Bugs identified and corrected.}\label{tab:fixes}\\
\toprule
\textbf{\#} & \textbf{Location} & \textbf{Original (buggy) behaviour} & \textbf{Corrected behaviour} \\
\midrule
\endfirsthead
\toprule
\textbf{\#} & \textbf{Location} & \textbf{Original} & \textbf{Corrected} \\
\midrule
\endhead
\bottomrule
\endlastfoot
1 & \texttt{WorkflowConfig}
  & Wrong \texttt{AGG\_FACTOR}; confusing two-step factor chain
  & Single \texttt{AGG\_FACTOR = 25} (50\,m / 2\,m); intermediate factors removed \\[4pt]
2 & BLOCK\,1
  & Terrain tools ran on native NAD83 bathymetry (EPSG:26920)
  & Reproject bathymetry to EPSG:32620 first; all derivatives inherit correct CRS \\[4pt]
3 & BLOCK\,3
  & Two-step chain (Slope $\to$ Copy $\to$ BlockStatistics) left output at 2\,m
  & Single \texttt{Aggregate(factor=25, MEAN)} on \texttt{slope\_2m.tif} $\Rightarrow$ true 50\,m cells \\[4pt]
4 & BLOCK\,4
  & \texttt{NbrCircle(5)} on 2\,m raster $=$ 10\,m neighbourhood
  & \texttt{NbrCircle(120)} on 2\,m raster $=$ 240\,m neighbourhood \\[4pt]
5 & BLOCK\,5
  & Slope-of-slope computed from a 2\,m raster mislabelled as ``50\,m''
  & Uses the true 50\,m slope produced by Fix~3 \\[4pt]
6 & BLOCK 6/6b
  & Arithmetic mean of aspect degrees — wrong at the 0°/360° wrap
  & Circular mean via sin/cos decomposition $\to$ aggregate $\to$ \texttt{atan2} recovery \\[4pt]
7 & BLOCK 7/7b
  & Only \texttt{curvature\_2m.tif} produced; aggregation deferred to R
  & \texttt{Aggregate(factor=25, MEAN)} added in Python; 50\,m output produced directly \\
\end{longtable}

\subsection{Mathematical Detail: Circular Mean for Aspect}

Aspect is a \textit{circular} variable (degrees, wrapping at 0°/360°).
A naive arithmetic mean fails near the wrap-around:
\[
  \overline{\theta}_{\text{naive}}(1°, 359°) = 180°
  \quad \text{(wrong; correct answer is } 0°)
\]
The circular mean is computed as follows:
\begin{align}
  \bar{s} &= \frac{1}{n}\sum_{i=1}^{n}\sin(\theta_i),\quad
  \bar{c}  = \frac{1}{n}\sum_{i=1}^{n}\cos(\theta_i) \\[4pt]
  \bar{\theta} &= \operatorname{atan2}(\bar{s},\,\bar{c})
                  \;\mathrm{mod}\; 2\pi
                  \quad\text{(converted to degrees)}
\end{align}
This is implemented in both the Python script (BLOCK\,6b) and the R helper
\texttt{circular\_mean\_aspect()}.

% =============================================================================
\section{Final Outputs and Verification}
% =============================================================================

BLOCK\,8 performs final verification.  All five outputs passed:

\begin{table}[h]
\centering
\caption{BLOCK 8 verification results.}
\rowcolors{2}{lightgray}{white}
\begin{tabular}{lllrr}
\toprule
\textbf{File} & \textbf{CRS} & \textbf{EPSG} & \textbf{Cell X\,(m)} & \textbf{Cell Y\,(m)} \\
\midrule
\cmark{} \texttt{slope\_50m.tif}           & WGS\_1984\_UTM\_Zone\_20N & 32620 & 50.0 & 50.0 \\
\cmark{} \texttt{slope\_240m\_50m.tif}     & WGS\_1984\_UTM\_Zone\_20N & 32620 & 50.0 & 50.0 \\
\cmark{} \texttt{slope\_of\_slope\_50m.tif}& WGS\_1984\_UTM\_Zone\_20N & 32620 & 50.0 & 50.0 \\
\cmark{} \texttt{aspect\_50m.tif}          & WGS\_1984\_UTM\_Zone\_20N & 32620 & 50.0 & 50.0 \\
\cmark{} \texttt{curvature\_50m.tif}       & WGS\_1984\_UTM\_Zone\_20N & 32620 & 50.0 & 50.0 \\
\bottomrule
\end{tabular}
\end{table}

\subsection{Intermediate 2\,m Archives}

\begin{table}[h]
\centering
\caption{Intermediate files kept for QA/archive.}
\rowcolors{2}{lightgray}{white}
\begin{tabularx}{\linewidth}{L{7cm}X}
\toprule
\textbf{Path (relative to \texttt{02\_terrain\_analysis\_outputs/})} & \textbf{Description} \\
\midrule
\texttt{00\_bathymetry\_source/STTSTJ\_2m\_native.tif}  & Original bathymetry in native CRS \\
\texttt{00\_bathymetry\_source/STTSTJ\_2m\_EPSG32620.tif} & Bathymetry reprojected to EPSG:32620 (2\,m) \\
\texttt{00\_intermediate/slope\_2m.tif}                 & Slope at native 2\,m resolution \\
\texttt{00\_intermediate/slope\_2m\_240m\_focal.tif}    & Focal-mean slope at 240\,m scale (2\,m pixels) \\
\texttt{03\_aspect/aspect\_2m.tif}                      & Aspect at 2\,m resolution \\
\texttt{02\_curvature/curvature\_2m.tif}                & Curvature at 2\,m resolution \\
\bottomrule
\end{tabularx}
\end{table}

% =============================================================================
\section{About the Agent and Model Used}
% =============================================================================

\begin{table}[h]
\centering
\caption{AI agent and model provenance.}
\rowcolors{2}{lightgray}{white}
\begin{tabularx}{\linewidth}{L{4cm}X}
\toprule
\textbf{Property} & \textbf{Value} \\
\midrule
Platform        & GitHub Copilot (\href{https://github.com}{github.com}) \\
Agent type      & Copilot Coding Agent / Task Agent \\
AI model        & Claude Sonnet (Anthropic) — version active in GitHub Copilot, April 2026 \\
Interface       & Asynchronous task-based: user submits task; agent works autonomously on a branch and reports progress via Pull Requests \\
Repository integration & Agent reads, creates, and edits files directly; all changes committed to a feature branch \\
Execution environment & Sandboxed Linux container; ArcGIS Pro not available — agent reasoned from code and user-reported errors \\
\bottomrule
\end{tabularx}
\end{table}

\textbf{Important note:} The AI agent \textit{did not execute the ArcPy code
itself}.  ArcGIS Pro is not available in the sandboxed cloud environment.
The agent reasoned about correctness from the code text, the error messages
reported by the user, and the unexpected output file sizes.  All actual raster
processing was performed by the user on their local ArcGIS Pro installation.

% =============================================================================
\section{How to Run the Code from Scratch}
\textit{(Step-by-step guide for non-experts)}
% =============================================================================

\subsection{Required Software}

\begin{table}[h]
\centering
\caption{Software prerequisites.}
\rowcolors{2}{lightgray}{white}
\begin{tabularx}{\linewidth}{L{3.5cm}L{5cm}X}
\toprule
\textbf{Software} & \textbf{Source} & \textbf{Notes} \\
\midrule
ArcGIS Pro 3.x  & Institution software portal  & Requires \textbf{Spatial Analyst Extension} \\
Jupyter Notebook & Installed with ArcGIS Pro or via \href{https://docs.conda.io}{conda} & Used to run blocks interactively \\
R $\geq$ 4.1    & \href{https://cran.r-project.org/}{cran.r-project.org} & Free; for exploratory R scripts \\
RStudio (optional) & \href{https://posit.co/}{posit.co} & Recommended R IDE \\
\bottomrule
\end{tabularx}
\end{table}

\subsection{Required Data}

\begin{itemize}[noitemsep]
  \item \textbf{\texttt{STTSTJ\_2m.tif}} — 2\,m bathymetry for St.\ Thomas /
        St.\ John.  Must be placed at the path defined in
        \texttt{WorkflowConfig.BATHYMETRY\_SOURCE\_TIF}.
  \item \textbf{WCMC coral-reef shapefile} — used only in the R exploratory
        script.
  \item \textbf{Hard-coral survey points} — \texttt{hard\_corals\_data\_filtered.shp},
        used only in the R exploratory script.
\end{itemize}

\subsection{Step 1 — Set Up Your Folder Structure}

The script creates output folders automatically, but understanding the layout
is helpful:

\begin{lstlisting}[style=shell]
02_terrain_analysis_outputs/
├── 00_bathymetry_source/   <- input + reprojected bathymetry
├── 00_intermediate/        <- temporary 2m files
├── 01_slope/               <- final slope outputs (50m)
├── 02_curvature/           <- curvature outputs (50m)
└── 03_aspect/              <- aspect outputs (50m)
\end{lstlisting}

\subsection{Step 2 — Open the Script in Jupyter Notebook}

\begin{enumerate}[noitemsep]
  \item Open \textbf{ArcGIS Pro}.
  \item In the ribbon: \textbf{Analysis $\to$ Python} to open the Python Notebook panel,
        or open Jupyter Notebook from the Start menu (the ArcGIS Pro \texttt{conda}
        environment includes Jupyter).
  \item Navigate to the folder containing
        \texttt{Coral\_Modeling\_Terrain\_Analysis\_USVI\_202604.py}.
  \item Open the file.  If using a \texttt{.py} extension, copy its contents
        into a new Jupyter Notebook (\texttt{.ipynb}) — one cell per BLOCK section.
\end{enumerate}

\subsection{Step 3 — Edit the File Paths}

Before running any block, update the two path variables to match your system:

\begin{lstlisting}[style=pycode]
# In WorkflowConfig class (top of the script):
PROJECT_ROOT = r"G:\Shared drives\NSF CoPE internal\..."  # <-- YOUR path
BATHYMETRY_SOURCE_TIF = r"G:\...\STTSTJ_2m.tif"          # <-- YOUR path
\end{lstlisting}

Use the \textbf{full path} to the folder / file on your computer.
The leading \texttt{r} before the quote means ``raw string'' and ensures
Windows backslashes are interpreted correctly.

\subsection{Step 4 — Run the Blocks in Order}

Run each BLOCK cell one at a time.  \textbf{Do not skip blocks} on the first
run.  Wait for each block to finish before running the next.

\begin{table}[h]
\centering
\caption{Expected runtime per BLOCK.}
\rowcolors{2}{lightgray}{white}
\begin{tabular}{llr}
\toprule
\textbf{BLOCK} & \textbf{Action} & \textbf{Time} \\
\midrule
BLOCK 0     & Load configuration, check extensions         & $<$5\,s \\
BLOCK 1     & Create folders; reproject bathymetry         & 1–2\,min \\
BLOCK 2     & Validate bathymetry file                     & $<$5\,s \\
Verification 2A & Print depth stats                         & $<$5\,s \\
BLOCK 3     & Slope at 2\,m, aggregate to 50\,m           & 1–2\,min \\
BLOCK 4     & Slope at 240\,m scale (large focal window)  & \textbf{5–8\,min} \\
BLOCK 5     & Slope-of-slope from true 50\,m slope        & 30–40\,s \\
BLOCK 6+6b  & Aspect at 2\,m; circular-mean aggregate to 50\,m & 2–3\,min \\
BLOCK 7+7b  & Curvature at 2\,m; aggregate to 50\,m       & 1–2\,min \\
BLOCK 8     & Final verification table                    & $<$10\,s \\
\bottomrule
\end{tabular}
\end{table}

Each block prints messages beginning with \textbf{\cmark} (success) or
\textbf{\xmark} (error).  \textit{Do not proceed to the next block if you see
an error message.}

\subsection{Step 5 — Verify the Final Outputs}

After BLOCK\,8 runs successfully, the console should display all five rows with
\cmark{} marks and cell size 50.0\,m (see Table\,5 above).

\subsection{Step 6 (Optional) — Run the R Exploratory Script}

\begin{enumerate}[noitemsep]
  \item Open \textbf{RStudio}.
  \item Open \texttt{01\_exploratory\_spatial\_analysis.R}.
  \item Edit the \texttt{root} variable at the top to match your folder path.
  \item Install required packages (first time only):
\end{enumerate}

\begin{lstlisting}[style=Rcode]
install.packages(c("terra", "sf", "dplyr"))
\end{lstlisting}

\begin{enumerate}[noitemsep, resume]
  \item Run the entire script (\texttt{Ctrl+Shift+Enter} in RStudio).
\end{enumerate}

The script will (a) clip the WCMC coral-reef polygon to the study area,
(b) report CRS and resolution of all layers, and (c) count survey points
inside/outside the coral mask.

\subsection{Troubleshooting Common Errors}

\begin{table}[h]
\centering
\caption{Common errors and solutions.}
\rowcolors{2}{lightgray}{white}
\begin{tabularx}{\linewidth}{L{5cm}L{4.5cm}X}
\toprule
\textbf{Error message} & \textbf{Likely cause} & \textbf{Solution} \\
\midrule
\texttt{FileNotFoundError: STTSTJ\_2m.tif not found}
  & Path in \texttt{WorkflowConfig} is wrong
  & Double-check the path; confirm the file exists \\
\texttt{Spatial Analyst Extension unavailable}
  & Extension not licensed / checked out
  & ArcGIS Pro $\to$ Project $\to$ Licensing $\to$ Check out Spatial Analyst \\
\texttt{PermissionError} writing to Google Drive
  & Drive sync conflict
  & Save to a local folder first; then copy to Drive \\
BLOCK\,4 finishes in $<$10\,s
  & Focal radius is wrong
  & Confirm \texttt{FOCAL\_RADIUS\_240M\_at2m = 120} \\
Output cell size is 2.0\,m instead of 50.0\,m
  & Wrong \texttt{AGG\_FACTOR}
  & Confirm \texttt{AGG\_FACTOR = 25} \\
\bottomrule
\end{tabularx}
\end{table}

% =============================================================================
\section{Data Sources and File Naming}
% =============================================================================

\subsection{Primary Input Data}

\begin{table}[h]
\centering
\caption{Input datasets.}
\rowcolors{2}{lightgray}{white}
\begin{tabularx}{\linewidth}{L{4cm}L{3.5cm}L{3cm}L{2cm}}
\toprule
\textbf{Dataset} & \textbf{Source} & \textbf{Native CRS} & \textbf{Resolution} \\
\midrule
\texttt{STTSTJ\_2m.tif} & NOAA / NSF CoPE & EPSG:26920 or :32620 & 2\,m \\
WCMC CoralReef 2021 & UNEP-WCMC & Geographic WGS84 & vector \\
Hard-coral survey points & NSF CoPE field surveys & Varies & vector \\
\bottomrule
\end{tabularx}
\end{table}

\subsection{Output Naming Convention}

\begin{itemize}[noitemsep]
  \item \texttt{*\_2m.tif} — computed at native 2\,m resolution (intermediate or archive).
  \item \texttt{*\_50m.tif} — final product; every cell is a true
        50\,m $\times$ 50\,m pixel in EPSG:32620.
  \item \texttt{*\_240m\_50m.tif} — slope smoothed over a 240\,m neighbourhood,
        aggregated to 50\,m cells.
\end{itemize}

% =============================================================================
\section{Glossary}
% =============================================================================

\begin{table}[h]
\centering
\rowcolors{2}{lightgray}{white}
\begin{tabularx}{\linewidth}{L{3.5cm}X}
\toprule
\textbf{Term} & \textbf{Definition} \\
\midrule
EPSG code
  & Numerical identifier for a coordinate reference system.
    EPSG:32620 = WGS\,84 / UTM Zone\,20N (used for most of the Caribbean). \\
CRS
  & Coordinate Reference System — the mathematical framework relating map
    coordinates to positions on Earth. \\
UTM
  & Universal Transverse Mercator — a projected CRS using metres as units. \\
Resolution / cell size
  & Side length of each raster pixel.
    50\,m means each pixel covers a 50\,m $\times$ 50\,m area on the ground. \\
Aggregation
  & Combining multiple small pixels into one larger pixel
    (e.g.\ $25 \times 25 = 625$ pixels at 2\,m merged into one pixel at 50\,m). \\
Focal statistics
  & Moving-window operation where each pixel is replaced by a summary statistic
    (e.g.\ mean) computed over a neighbourhood. \\
Aspect
  & Compass direction (0°–360°) that a slope face points toward. \\
Curvature
  & Rate of change of slope.
    Positive $=$ convex (ridgeline); Negative $=$ concave (valley). \\
Circular mean
  & Correct method for averaging angular measurements that wrap at 0°/360°. \\
\texttt{arcpy}
  & Python package bundled with ArcGIS Pro providing access to all ArcGIS
    geoprocessing tools. \\
\texttt{terra}
  & R package for raster and vector spatial analysis. \\
Snap raster
  & Reference raster forcing all outputs to align to the same grid origin and
    cell boundaries. \\
\bottomrule
\end{tabularx}
\end{table}

% =============================================================================
%  FOOTER
% =============================================================================

\vfill
\begin{center}
  \color{darkgray}\small
  \rule{0.5\linewidth}{0.4pt}\\[4pt]
  This document was generated by the \textbf{GitHub Copilot Coding Agent}
  (Claude Sonnet, Anthropic) in collaboration with \textbf{J.\ Olaya},
  April 2026.\\
  It should be updated if the code or workflow changes significantly.\\[4pt]
  Repository:
  \href{https://github.com/jolaya80/jolaya80}{\texttt{github.com/jolaya80/jolaya80}}
\end{center}

\end{document}
